\section{Evaluation Methods for Code Translation}

% Nem értem... Az előbb már írtál a benchmarkokról, de itt is írsz a kiértékelésről... Kicsit nem áll össze a kép, át kéne gondolni

Evaluating the quality of code translation remains a challenging task, as it
requires assessing not only syntactic correctness but also the preservation of
program semantics. Existing evaluation approaches can be broadly divided into
text-based metrics and execution-based methods.

Text-based metrics are derived from natural language processing and measure the
similarity between generated and reference code at the token level. One of the
most widely used metrics is BLEU, which evaluates n-gram overlap between the
generated output and the reference translation \cite{bleu_metric}. Although BLEU
is simple and widely adopted, it is limited in the context of code, as it does
not account for syntactic structure or semantic equivalence.

To address these limitations, CodeBLEU was introduced as an extension of BLEU
specifically designed for code evaluation. CodeBLEU incorporates additional
factors such as abstract syntax tree (AST) structure and data-flow information,
providing a more comprehensive assessment of code quality
\cite{ren2020codebleu}. However, despite these improvements, CodeBLEU still
relies on surface-level similarity and may fail to capture functional
correctness.

Execution-based evaluation methods provide a more reliable alternative by
assessing whether the generated code behaves correctly when executed. These
approaches typically involve compiling the translated code and running it
against a set of test cases. The correctness of the output is then compared to
the expected results, often using metrics such as pass rate or pass@k. This
method directly measures whether the translated code preserves the intended
functionality.

Recent studies have shown that models achieving high scores on text-based
metrics do not necessarily perform well in execution-based evaluations,
highlighting the limitations of similarity-based approaches
\cite{huang2022exeds}. In particular, code with low textual similarity may still
be functionally correct, while code with high similarity may fail to execute
correctly.

Overall, the discrepancy between text-based and execution-based evaluation
highlights a fundamental challenge in code translation: semantic correctness
cannot be reliably captured by token-level similarity alone. This motivates the
increasing use of execution-based evaluation methods, which better reflect
real-world correctness and are therefore more suitable for assessing practical
code translation systems.